%% kapa: exam=1 round=24 year=2013 period=2 form=A subject=accounting subject_ko=회계학 seq=16 number=16
%% source: Q-Net 감정평가사 기출문제 (cst00309) artlSeq=5073802 file=extracted/1008/5073802/2013년 제24회 감정평가사 1차시험 기출문제/2교시(A형).hwp
%% conversion: kordoc 4.24.51 -> tools/gen1.py (auto); review=proofread-sonnet; flags=none
\begin{question}{16}
(주)감평은 20×1년 4월 1일에 다음과 같은 조건으로 발행된 채무증권을 취득하고 매도가능금융자산으로 분류하였다.
\begin{framedtext}
○ 액면금액: ￦1,000,000 \par
○ 발 행 일: 20×1년 4월 1일 \par
○ 상 환 일: 20×4년 3월 31일 \par
○ 표시이자율: 연 8\% \par
○ 유효이자율: 연 10\% \par
○ 이자는 매년 3월 31일에 지급하고 월할계산한다. \par
○ 20×1년 12월 31일 현재 공정가치(기간 경과이자 포함)는 ￦1,010,000이다.
\end{framedtext}
20×1년 12월 31일에 (주)감평이 인식해야 할 매도가능금융자산평가손익은? (단, 현가계수는 아래의 표를 이용하며, 단수차이로 인한 오차가 있으면 가장 근사치를 선택한다.) \par
\uline{현가계수표} \par
{\small\centering
\begin{tabular}{|c|c|c|}
\hline
\makecell{(3년 기준)} & \makecell{8\%} & \makecell{10\%} \\
\hline
\makecell{단일금액 ￦1의 현가계수} & \makecell{0.7938} & \makecell{0.7513} \\
\hline
\makecell{정상연금 ￦1의 현가계수} & \makecell{2.5771} & \makecell{2.4868} \\
\hline
\end{tabular}\par}

\begin{choices}
  \item 이익 ￦9,091
  \item 이익 ￦48,488
  \item 손실 ￦10,244
  \item 손실 ￦11,268
  \item 손실 ￦11,512
\end{choices}
\end{question}
\answer{5}
